\chapter{C1-Esterase Inhibitor}

\section*{What Is This Test?}
C1-esterase inhibitor testing measures the concentration and functional activity of a complement-regulating protein. It is used in the evaluation of hereditary or acquired angioedema, conditions that can cause recurrent swelling without typical hives.

\section*{Alternative Names}
C1-INH; C1 Inhibitor Antigen; C1 Inhibitor Functional Assay; Hereditary Angioedema Testing.

\section*{Why This Test Is Ordered}
Testing is appropriate for recurrent unexplained swelling of the face, lips, limbs, bowel, or airway, especially when episodes are not accompanied by urticaria or itching. C4 is usually measured with C1-inhibitor concentration and function.

\section*{How This Test Is Performed}
A blood sample is collected and analyzed for C1-inhibitor protein concentration and functional activity. Complement C4 and, in selected cases, C1q help distinguish hereditary from acquired forms.

\section*{Preparation, Risks, and Limitations}
No fasting is required. Venipuncture is the usual risk. Complement proteins can be consumed during an acute attack, and delayed processing or poor temperature control can invalidate functional testing. Results should be repeated or confirmed when they conflict with the clinical history.

\section*{Understanding Results}
Low C1-inhibitor concentration and/or function with low C4 supports C1-inhibitor deficiency. Acquired disease may have low C1q, while hereditary disease often has normal C1q. Normal results do not exclude all forms of angioedema.

\section*{References}
World Allergy Organization and European Academy of Allergy and Clinical Immunology guidance on hereditary angioedema.

\clearpage
\section*{Understanding Results and Follow-up}
The laboratory report should identify the specimen, method, units, reference interval or decision limit, and whether the result is positive, negative, abnormal, or inconclusive. Reference intervals vary with age, sex, pregnancy, specimen type, assay, and laboratory; a result outside a range is not automatically a diagnosis, and a result within a range does not always exclude disease. Discuss unexpected results with the ordering clinician, who can decide whether repeat or confirmatory testing, treatment, or referral is appropriate. Do not change medicines or delay urgent care based on an isolated result.


\section*{Regulatory Status and Authorization Date}
Regulatory status applies to the specific assay, instrument, manufacturer, and intended use rather than to the test name alone. No single approval date can be assigned to this test category. For a commercial assay, verify the current product in the relevant regulator's database: the U.S. Food and Drug Administration (FDA) clearance, approval, or authorization; India's Central Drugs Standard Control Organisation (CDSCO) licence or permission; the European Union In Vitro Diagnostic Regulation (EU IVDR) CE certificate issued through the applicable conformity-assessment route; or the United Kingdom Medicines and Healthcare products Regulatory Agency (MHRA) registration and UKCA/CE status. Laboratory-developed tests may not have an individual FDA, CDSCO, EU IVDR, or MHRA approval date.
\clearpage
